\begin{document}
\lhead[ \leftmark ]{Informatik, Interessenwoche ``Web-Technologien''}
\rhead[Informatik]{Cyril Wendl \& Thomas Graf, \today, Winterthur}

\lstset{language=HTML}
\section*{Auftrag}

Im Rahmen dieses Blockkurses erstellen Sie Ihre eigene Webseite über ein Thema rund um das Web. Dabei lernen Sie insbesondere die Themen HTML, CSS und JavaScript kennen. Wählen Sie sich dazu alleine oder in 2er- bis 3er-Gruppen ein für Sie spannendes Thema zum Thema Internet / Informatik aus (s. Artikel-Auswahl auf Moodle für eine Themenübersicht).

\subsection*{Anforderungen}

Ihre Webseite sollte dieses Thema auf schöne und einfache Art zusammenfassen, in dem sie folgende Elemente verwendet:
\begin{todolist}
    \item \href{https://www.w3schools.com/html/html_basic.asp}{Einfache HTML-Elemente} wie Überschriften (\lstinline|<h1>...</h1>|), Absätze (\lstinline|<p>...</p>|) und Zeilenumbrüche (\lstinline|<br/>|)
    \item Ein sinnvoller \href{https://www.w3schools.com/html/html_page_title.asp}{Seiten-Titel} mit dem \lstinline|<title>...</title>|-Tag
    \item \href{https://www.w3schools.com/html/html_formatting.asp}{Formatierungs-Tags} wie etwa \lstinline|<strong>...</strong>, <em>...</em>| usw.
    \item \href{https://www.w3schools.com/html/html_images.asp}{Bilder}, mit dem Tag \lstinline|<img src="...">...</img>|
    \item Eine oder mehrere \href{https://www.w3schools.com/html/html_tables.asp}{Tabellen} mit \lstinline|<table>...</table>|, \lstinline|<tr>...</tr>|, \lstinline|<th>...</th>| und \lstinline|<td>...</td>|
    \item Eine oder mehrere \href{https://www.w3schools.com/html/html_lists.asp}{geordnete und ungeordnete Listen} mit \lstinline|<ul>...</ul>|, \lstinline|<ol>...</ol>| und \lstinline|<li>...</li>|
    \item Ein oder mehrere passende \href{https://www.w3schools.com/html/html_quotation_elements.asp}{Zitate}
    \item \href{https://www.w3schools.com/html/html_links.asp}{Links}  zu weiterführenden Webseiten / Quellen, mit dem Tag \lstinline|<a href="...">...</a>|
    \item Mindestens 3 weitere HTML-Tags, die Sie dem Tutorial auf \href{https://www.w3schools.com/html/html_quotation_elements.asp}{w3schools} entnehmen können.
    \item \href{https://www.w3schools.com/html/html_styles.asp}{Stil-Elemente direkt in einem HTML-Tag} mit dem \lstinline|<... style="...">|-Attribut
    \item Verwendung eines \href{https://www.w3schools.com/html/html_css.asp}{externen Stylesheets}
    \item Validierung: Der \href{https://validator.w3.org/nu/#textarea}{HTML}- sowie der \href{https://jigsaw.w3.org/css-validator/#validate_by_input}{CSS}-Code sollen fehlerfrei sein.
\end{todolist}

\subsection*{Bonus}

\begin{todolist}
    \item Wenn Ihre Webseite ein eigenes \href{https://www.w3schools.com/html/html_favicon.asp}{Favicon} verwendet
    \item \href{https://www.w3schools.com/html/html_basic.asp}{Interaktivität mit JavaScript} oder \href{https://www.w3schools.com/jquery/default.asp}{jQuery}, beispielsweise um weitere Informationen interaktiv anzuzeigen, Popups zu erstellen usw.
    \item Vektor-Grafiken in Form von SVG-Dateien, beispielsweise um Ihr eigenes Logo darzustellen
\end{todolist}

\subsection*{Abgabe der Webseite}
Wir werden am Schluss alle Webseiten auf eine öffentliche Web-Adresse hochladen. Reichen Sie dazu bitte ihre Dateien via \href{https://moodle.ksimlee.ch/mod/assign/view.php?id=32219}{Moodle} ein.

\subsection*{Prämierung der besten Webseiten}


\emoji{trophy} Die besten 3 Webseiten werden mit einer kleinen Auszeichnung und einem Preis belohnt!
\end{document}